\documentclass[11pt]{article}

\usepackage[margin=1in]{geometry}
\usepackage{amsmath,amssymb,amsthm}
\usepackage{enumitem}
\usepackage{hyperref}
\usepackage{xcolor}
\usepackage{booktabs}

\hypersetup{colorlinks=true, linkcolor=blue, urlcolor=blue, citecolor=blue}

\title{TOAE209/TOAH209 -- Design and Analysis of Algorithms\\[4pt]
\large Week 9: Theoretical Questions}
\author{Foreign Trade University}
\date{}

\begin{document}
\maketitle

\section*{Instructions}

Part A contains 16 questions requiring written explanations or short proofs. Part B is a
10-question quiz (true/false and multiple choice). Attempt every question on your own
before consulting Part C (Solutions) at the end of this document.

\section*{Part A: Theory Questions}

\begin{enumerate}[label=\textbf{A\arabic*.}, leftmargin=2.2em]

\item State the four-step ``recipe'' for designing a dynamic-programming algorithm
(subproblems, recurrence, evaluation order, reconstruction). Which step is usually the
hardest, and why?

\item State the Sequence Alignment problem precisely: what are the inputs (the two strings,
the gap cost $\delta$, and the mismatch costs $\alpha_{pq}$), what is a valid alignment, and
what is the objective? How does Levenshtein edit distance arise as a special case?

\item Derive the alignment recurrence by reasoning about the \emph{last column} of an
optimal alignment of $x_1\cdots x_i$ and $y_1\cdots y_j$. List the three cases and state the
base cases.

\item Prove that the alignment recurrence is correct (i.e.\ that $\mathrm{OPT}(i,j)$ equals
the true minimum-cost alignment of the two prefixes). Make clear where you use the fact that
the rest of an optimal alignment, after fixing the last column, must itself be optimal.

\item What are the time and space complexities of the basic table-filling alignment
algorithm? Explain why the time is $O(mn)$ in terms of (number of subproblems) $\times$
(work per subproblem).

\item Describe how to reconstruct an actual optimal alignment (the two padded strings with
gap symbols) once the table is filled. Why does the traceback take only $O(m+n)$ time even
though the table has $O(mn)$ cells?

\item Explain, at a high level, Hirschberg's idea for computing an optimal alignment in only
$O(m+n)$ space (while keeping $O(mn)$ time). What single observation about the dependency
structure of the table makes the space saving possible, and what extra idea
(divide-and-conquer) is needed to still recover the alignment, not just its cost?

\item State the Longest Common Subsequence (LCS) recurrence and explain the reasoning behind
each case (when $x_i = y_j$ versus $x_i \ne y_j$). In what sense is LCS a special case of the
sequence-alignment framework?

\item Prove the identity $d(X,Y) = m + n - 2\,\mathrm{LCS}(X,Y)$, where $d$ is the minimum
number of single-character insertions and deletions (no substitutions) needed to turn $X$
(length $m$) into $Y$ (length $n$).

\item Explain why Dijkstra's algorithm can return incorrect answers on graphs with negative
edge weights. Give (or describe) a small example where Dijkstra fails. Also explain why
``shortest path'' is ill-defined when a negative cycle is reachable.

\item State the Bellman--Ford DP: define the subproblem $\mathrm{OPT}(k,v)$, write the
recurrence (reasoning about the last edge), and state the base cases. Why is $k$ bounded by
$n-1$ when there is no negative cycle?

\item Prove that after $n-1$ passes of edge relaxation, Bellman--Ford computes correct
shortest-path distances in a graph with no negative cycle. State the running time and
justify it.

\item Describe how Bellman--Ford detects a negative cycle, and how one can recover the actual
cycle (not just a yes/no flag). Why does running exactly one additional ($n$-th) relaxation
pass suffice as the test?

\item State the Floyd--Warshall recurrence: define $d_k(i,j)$ in terms of allowed
intermediate vertices, write the recurrence (reasoning about whether vertex $k$ is used),
and state the base case. What is the running time, and how does a negative cycle reveal
itself in the final matrix?

\item Explain why shortest (and even \emph{longest}) paths in a DAG can be computed in
$O(n+m)$ time by relaxing edges in topological order. Why is the longest-path problem
tractable on a DAG but NP-hard on general graphs?

\item State the Matrix Chain Multiplication recurrence (define $M(i,j)$, the split-point
minimisation, and the base case) and the Grid Minimum-Path-Sum recurrence. For each, give
the number of subproblems and the work per subproblem, and hence the total running time.

\end{enumerate}

\newpage
\section*{Part B: Quiz}

\begin{enumerate}[label=\textbf{B\arabic*.}, leftmargin=2.2em]

\item \textbf{(Multiple choice)} The sequence-alignment dynamic program runs in time:
\begin{enumerate}[label=(\alph*)]
    \item $O(m + n)$
    \item $O(m n)$
    \item $O(2^{m+n})$
    \item $O(mn \log(mn))$
\end{enumerate}

\item \textbf{(True/False)} In the alignment recurrence, the cell $\mathrm{OPT}(i,j)$
depends only on the three cells $\mathrm{OPT}(i-1,j-1)$, $\mathrm{OPT}(i-1,j)$, and
$\mathrm{OPT}(i,j-1)$.

\item \textbf{(Short answer)} The Levenshtein edit distance between \texttt{kitten} and
\texttt{sitting} is what number, and what three single-character operations achieve it?

\item \textbf{(Multiple choice)} For strings of lengths $m$ and $n$, the relation between
insert/delete-only edit distance $d$ and LCS length $L$ is:
\begin{enumerate}[label=(\alph*)]
    \item $d = m + n - L$
    \item $d = m + n - 2L$
    \item $d = \max(m,n) - L$
    \item $d = 2L - m - n$
\end{enumerate}

\item \textbf{(True/False)} Dijkstra's algorithm always computes correct shortest paths,
even when some edge weights are negative, as long as there is no negative cycle.

\item \textbf{(Multiple choice)} Bellman--Ford runs in time:
\begin{enumerate}[label=(\alph*)]
    \item $O(m \log n)$
    \item $O(m + n)$
    \item $O(mn)$
    \item $O(n^3)$
\end{enumerate}

\item \textbf{(Short answer)} After running Bellman--Ford's $n-1$ relaxation passes, you run
one more pass and find an edge $(u,v,w)$ with $\mathrm{dist}[u] + w < \mathrm{dist}[v]$.
What does this tell you about the graph?

\item \textbf{(Multiple choice)} Floyd--Warshall computes all-pairs shortest paths in time:
\begin{enumerate}[label=(\alph*)]
    \item $O(n^2)$
    \item $O(n^3)$
    \item $O(mn)$
    \item $O(n^2 \log n)$
\end{enumerate}

\item \textbf{(True/False)} In a directed acyclic graph, both shortest \emph{and} longest
paths from a source can be computed in $O(n+m)$ time by relaxing edges in topological order.

\item \textbf{(Short answer)} In matrix chain multiplication, $M(i,j)$ is computed by
minimising over split points $k$ with $i \le k < j$. Write the cost contributed by choosing
split point $k$ (in terms of $M$ and the dimensions $p_{i-1}, p_k, p_j$).

\end{enumerate}

\newpage
\section*{Part C: Solutions}

\subsection*{Part A}

\begin{enumerate}[label=\textbf{A\arabic*.}, leftmargin=2.2em]

\item The recipe: \textbf{(1) Define subproblems} -- a small family of subproblems, each
indexed by one or more parameters, whose optima combine into the whole. \textbf{(2) Write a
recurrence} -- express a subproblem's optimum in terms of smaller subproblems' optima, by
reasoning about the \emph{last decision} an optimal solution makes. \textbf{(3) Order the
computation} -- evaluate subproblems so each is solved before it is needed (e.g.\ smallest
first, or in topological order). \textbf{(4) Reconstruct} -- trace back (or store parent /
choice pointers) to recover an actual optimal solution, not just its value. Step (2) is
usually hardest: identifying the right notion of ``last decision'' and proving the rest of
an optimal solution decomposes into independent, already-defined subproblems. Steps (1),
(3), (4) are then largely mechanical, and the running time is (number of subproblems)
$\times$ (work per subproblem).

\item \textbf{Input:} two strings $X = x_1\cdots x_m$, $Y = y_1\cdots y_n$; a gap cost
$\delta \ge 0$; and mismatch costs $\alpha_{pq} \ge 0$ with $\alpha_{pp} = 0$. A
\textbf{valid alignment} writes the two strings one above the other with gap symbols
inserted so both have equal length and no column is two gaps; each column is a match
(equal characters, cost $0$), a mismatch (cost $\alpha_{x_i y_j}$), or a character against a
gap (cost $\delta$). The \textbf{objective} is an alignment of minimum total column-cost
sum. \textbf{Levenshtein edit distance} is the special case $\delta = 1$, $\alpha_{pq} = 1$
for $p \ne q$: it counts the minimum number of insertions, deletions, and substitutions to
transform $X$ into $Y$.

\item Look at the last column of an optimal alignment of $x_1\cdots x_i$ and $y_1\cdots
y_j$. Order is preserved, so exactly one of three things happens to $x_i$ and $y_j$:
\emph{(a)} they are aligned together -- cost $\alpha_{x_i y_j}$ plus an optimal alignment of
$x_1\cdots x_{i-1}$ and $y_1\cdots y_{j-1}$; \emph{(b)} $x_i$ is aligned with a gap (deleted)
-- cost $\delta$ plus an optimal alignment of $x_1\cdots x_{i-1}$ and $y_1\cdots y_j$;
\emph{(c)} $y_j$ is aligned with a gap (inserted) -- cost $\delta$ plus an optimal alignment
of $x_1\cdots x_i$ and $y_1\cdots y_{j-1}$. Hence
\[
\mathrm{OPT}(i,j) = \min\big( \alpha_{x_i y_j} + \mathrm{OPT}(i-1,j-1),\ \delta +
\mathrm{OPT}(i-1,j),\ \delta + \mathrm{OPT}(i,j-1)\big),
\]
with base cases $\mathrm{OPT}(i,0) = i\delta$ and $\mathrm{OPT}(0,j) = j\delta$.

\item By induction on $i+j$. \emph{Base cases:} aligning a non-empty prefix with the empty
string forces one gap per character, cost $i\delta$ (resp.\ $j\delta$); these are obviously
optimal (and the only) alignments. \emph{Inductive step} ($i,j\ge 1$): take any optimal
alignment $A$ of the two prefixes and look at its last column, which is one of the three
cases. Suppose it aligns $x_i$ with $y_j$. Removing that column leaves an alignment $A'$ of
$x_1\cdots x_{i-1}$ and $y_1\cdots y_{j-1}$, and $\mathrm{cost}(A) = \alpha_{x_i y_j} +
\mathrm{cost}(A')$. \emph{Here is the key step:} $A'$ must be an \emph{optimal} alignment of
those shorter prefixes -- otherwise replacing $A'$ by a cheaper alignment and re-appending
the last column would beat $A$, contradicting $A$'s optimality. By the inductive hypothesis
$\mathrm{cost}(A') = \mathrm{OPT}(i-1,j-1)$, so $\mathrm{cost}(A) = \alpha_{x_i y_j} +
\mathrm{OPT}(i-1,j-1)$. The deletion and insertion cases are symmetric. Thus
$\mathrm{cost}(A)$ equals one of the three expressions, so $\mathrm{OPT}(i,j) \ge \min(\cdots)
= \mathrm{cost}(A)$ would be wrong -- rather, $\mathrm{cost}(A) \ge \min(\cdots)$ since the
min is over all three options. Conversely each of the three expressions is realised by an
actual alignment (optimal alignment of the relevant smaller prefixes plus the matching last
column), so the minimum is achievable. Hence $\mathrm{OPT}(i,j) = \mathrm{cost}(A) =
\min(\cdots)$.

\item There are $(m+1)(n+1) = O(mn)$ subproblems (cells). Each cell is computed by taking
the minimum of three already-computed neighbouring cells plus $O(1)$ arithmetic, i.e.\
$O(1)$ work per subproblem. Total time $= O(mn) \times O(1) = O(mn)$. The basic algorithm
stores the whole table, so space is also $O(mn)$.

\item Trace back from cell $(m,n)$ toward $(0,0)$. At each cell $(i,j)$ determine which of
the three terms achieved $\mathrm{OPT}(i,j)$: if the diagonal term, emit a column pairing
$x_i$ with $y_j$ and move to $(i-1,j-1)$; if the ``up'' term, emit $x_i$ against a gap and
move to $(i-1,j)$; if the ``left'' term, emit a gap against $y_j$ and move to $(i,j-1)$.
Stop at $(0,0)$ and reverse the emitted columns. Each step decreases $i+j$ by at least $1$,
so the traceback visits at most $m+n$ cells -- hence $O(m+n)$ time -- even though filling the
table took $O(mn)$ (the traceback follows a single path through the table, not the whole
grid).

\item \textbf{Observation:} row $i$ of the table depends only on row $i-1$ (each cell uses
the cell above, to the left, and diagonally up-left). So computing just the final
\emph{cost} $\mathrm{OPT}(m,n)$ needs only two rows in memory, i.e.\ $O(\min(m,n))$ space.
But two rows do not retain enough information to trace the alignment back.
\textbf{Hirschberg's extra idea:} divide and conquer. An optimal alignment must pass through
some cell $(m/2, q)$ in the middle row. Compute, in linear space, the cost of optimally
aligning the top half of $X$ with each prefix $y_1\cdots y_q$ (a forward sweep) and the cost
of aligning the bottom half of $X$ with each suffix $y_{q+1}\cdots y_n$ (a backward sweep);
the optimal $q$ minimises the sum. Recurse on the two independent halves. The work satisfies
$T(m,n) = O(mn) + T(m/2, q) + T(m/2, n-q)$, which solves to $O(mn)$ time, while the space is
$O(m+n)$.

\item LCS recurrence: $L(i,j) = 0$ if $i=0$ or $j=0$; $L(i,j) = 1 + L(i-1,j-1)$ if
$x_i = y_j$; and $L(i,j) = \max(L(i-1,j), L(i,j-1))$ if $x_i \ne y_j$. \emph{Reasoning:} if
$x_i = y_j$ there is always an optimal LCS that matches them (matching equal last characters
never hurts), reducing to the LCS of the two shorter prefixes plus one. If $x_i \ne y_j$,
they cannot \emph{both} be the last matched character, so at least one is unused; dropping
$x_i$ gives $L(i-1,j)$ and dropping $y_j$ gives $L(i,j-1)$, and we take the better. LCS is
the special case of alignment that \emph{maximises} score with match reward $+1$, mismatch
$0$, and gap $0$ -- the score of an alignment is then exactly the number of matched
(equal) columns, i.e.\ the length of a common subsequence.

\item ($\le$) Let $S$ be a longest common subsequence, $|S| = L$. Build an edit script that
keeps the $L$ characters of $S$ in place: delete the $m - L$ characters of $X$ not in $S$,
then insert the $n - L$ characters of $Y$ not in $S$. This transforms $X$ into $Y$ using
$(m-L) + (n-L) = m + n - 2L$ insert/delete operations, so $d(X,Y) \le m + n - 2L$. ($\ge$)
Take any optimal insert/delete script turning $X$ into $Y$. The characters of $X$ that are
\emph{never deleted} survive into $Y$ in their original relative order, so they form a common
subsequence; say there are $\ell$ of them, with $\ell \le L$. The script deletes $m - \ell$
characters of $X$ and inserts $n - \ell$ characters to build $Y$, for $m + n - 2\ell \ge m +
n - 2L$ operations. Hence $d(X,Y) \ge m + n - 2L$. Combining, $d(X,Y) = m + n - 2L$.

\item Dijkstra finalises a vertex's distance the moment it is extracted from the priority
queue (the minimum-tentative-distance unvisited vertex), assuming no later edge can reduce
it -- valid only with non-negative weights. With a negative edge, a vertex finalised early
might later be reached more cheaply via a path through a negative edge, and Dijkstra never
revisits it. \emph{Example:} vertices $s, a, b$ with edges $s \to a$ (weight $1$),
$s \to b$ (weight $4$), $b \to a$ (weight $-3$). Dijkstra finalises $a$ at distance $1$
(extracted before $b$), but the true shortest distance to $a$ is $4 + (-3) = 1$ here -- to
make it fail use $s\to a = 2$, $s\to b = 1$, $b \to a = -2$: Dijkstra finalises $a=2$, but
$s\to b\to a = 1 - 2 = -1$ is shorter. A \textbf{negative cycle} reachable en route makes
``shortest path'' ill-defined: looping the cycle repeatedly lowers the cost without bound,
so the infimum is $-\infty$ and no finite shortest path exists.

\item \textbf{Subproblem:} $\mathrm{OPT}(k,v)$ = minimum cost of an $s$--$v$ path using at
most $k$ edges. \textbf{Recurrence} (last edge): either an optimal $\le k$-edge path uses
$\le k-1$ edges, or its last edge is some $(u,v)$ preceded by an optimal $\le k-1$-edge path
to $u$:
\[
\mathrm{OPT}(k,v) = \min\Big( \mathrm{OPT}(k-1, v),\ \min_{(u,v)\in E}\big(\mathrm{OPT}(k-1,u)
+ w(u,v)\big)\Big),
\]
with $\mathrm{OPT}(0,s)=0$ and $\mathrm{OPT}(0,v) = \infty$ for $v \ne s$. \textbf{Bound on
$k$:} if there is no negative cycle, every shortest path is \emph{simple} (repeating a vertex
would mean a cycle, and a non-negative cycle can be removed without increasing cost), so it
has at most $n-1$ edges; thus $\mathrm{OPT}(n-1, v)$ already gives the true distances.

\item By induction on the pass number $k$, the invariant ``after pass $k$,
$\mathrm{dist}[v] \le \mathrm{OPT}(k,v)$ for all $v$''. \emph{Base:} before any pass,
$\mathrm{dist}[s]=0=\mathrm{OPT}(0,s)$, all others $\infty$. \emph{Step:} assume it holds
after pass $k-1$. Fix $v$ and an optimal $\le k$-edge path $P$. If $P$ uses $\le k-1$ edges,
$\mathrm{dist}[v]$ is already $\le \mathrm{OPT}(k-1,v) = \mathrm{cost}(P)$. Otherwise $P$'s
last edge $(u,v)$ follows a $\le k-1$-edge prefix of cost $\mathrm{OPT}(k-1,u)$; by
hypothesis $\mathrm{dist}[u] \le \mathrm{OPT}(k-1,u)$ entering pass $k$, and relaxing
$(u,v)$ in pass $k$ sets $\mathrm{dist}[v] \le \mathrm{dist}[u] + w \le \mathrm{cost}(P)$.
With no negative cycle, shortest paths are simple ($\le n-1$ edges), so after $n-1$ passes
$\mathrm{dist}[v]$ equals the true distance. \emph{Running time:} each pass relaxes all $m$
edges and there are $n-1$ passes, so $O(mn)$.

\item \textbf{Detection:} after the $n-1$ passes, run one more pass; if any edge $(u,v,w)$
still satisfies $\mathrm{dist}[u] + w < \mathrm{dist}[v]$, report a negative cycle.
\textbf{Why one pass suffices:} in a negative-cycle-free graph, distances have converged
after $n-1$ passes (previous question), so \emph{no} edge can be further relaxed; therefore
a successful relaxation on the $n$-th pass means distances are still decreasing, which can
only happen if a reachable negative cycle lets paths grow beyond $n-1$ edges while still
improving. \textbf{Recovering the cycle:} when the $n$-th-pass relaxation of $(u,v)$
succeeds, set $\mathrm{parent}[v] = u$, then start at $v$ and follow parent pointers $n$
times (this guarantees ending at a vertex \emph{on} the cycle, since the parent chain must
eventually enter the cycle); call that vertex $x$, then follow parents from $x$ until $x$
recurs -- the vertices traversed form the negative cycle.

\item \textbf{Subproblem:} $d_k(i,j)$ = length of the shortest $i$--$j$ path whose
\emph{intermediate} vertices all lie in $\{1,\ldots,k\}$. \textbf{Recurrence} (is vertex
$k$ used?):
\[
d_k(i,j) = \min\big( d_{k-1}(i,j),\ d_{k-1}(i,k) + d_{k-1}(k,j)\big),
\]
with $d_0(i,j) = w(i,j)$ for edges, $d_0(i,i)=0$, else $\infty$. Either the optimal path
avoids $k$ (first term), or it passes through $k$ exactly once, splitting into an $i \to k$
and a $k \to j$ path each using only $\{1,\ldots,k-1\}$ as intermediates (second term). With
$n$ values of $k$ and $n^2$ pairs $(i,j)$, each $O(1)$, the time is $O(n^3)$ and space
$O(n^2)$ (one matrix updated in place). A \textbf{negative cycle} shows up as some diagonal
entry $d_n(i,i) < 0$ (a path from $i$ back to $i$ of negative length).

\item A DAG has a \textbf{topological order} in which every edge goes from an earlier to a
later vertex. Processing vertices in this order, when we reach $v$ all of its predecessors
$u$ (with $(u,v)\in E$) are already finalised, so $\mathrm{dist}(v) = \min_{(u,v)\in E}
(\mathrm{dist}(u) + w)$ (or $\max$ for longest paths) is computed in one sweep. Each edge is
examined exactly once, giving $O(n+m)$ (plus $O(n+m)$ for the topological sort). Longest
paths are tractable on a DAG precisely because there are no cycles -- the recurrence is
well-founded and longest \emph{walks} equal longest \emph{paths}. On a general graph, a
longest \emph{simple} path is NP-hard (e.g.\ a Hamiltonian path is a longest simple path of
$n-1$ edges), and longest walks are unbounded if any positive cycle exists, so no analogous
DP works.

\item \textbf{Matrix chain:} $M(i,j)$ = minimum scalar multiplications to compute
$A_i\cdots A_j$, with $A_t$ of dimension $p_{t-1}\times p_t$:
\[
M(i,j) = \min_{i\le k < j}\big( M(i,k) + M(k+1,j) + p_{i-1}p_k p_j\big), \qquad M(i,i)=0.
\]
There are $O(n^2)$ subproblems $(i,j)$, each minimising over $O(n)$ split points, so
$O(n^3)$ time. \textbf{Grid minimum path sum:} with cell costs $c(i,j)$ and moves
right/down,
\[
S(i,j) = c(i,j) + \min\big( S(i-1,j),\ S(i,j-1)\big),
\]
($S(0,0) = c(0,0)$, first row/column accumulate). There are $O(rc)$ subproblems, each $O(1)$
work, so $O(rc)$ time.

\end{enumerate}

\subsection*{Part B}

\begin{enumerate}[label=\textbf{B\arabic*.}, leftmargin=2.2em]

\item \textbf{(b) $O(mn)$.} There are $(m+1)(n+1)$ cells, each computed in $O(1)$.

\item \textbf{True.} Each alignment cell uses the cell diagonally up-left (match/mismatch),
directly above (deletion), and directly left (insertion).

\item Edit distance $\mathbf{3}$. One optimal sequence: substitute \texttt{k}$\to$\texttt{s}
(\texttt{kitten}$\to$\texttt{sitten}), substitute \texttt{e}$\to$\texttt{i}
(\texttt{sitten}$\to$\texttt{sittin}), insert \texttt{g} at the end
(\texttt{sittin}$\to$\texttt{sitting}).

\item \textbf{(b) $d = m + n - 2L$.} Each unmatched character of $X$ is deleted and each
unmatched character of $Y$ is inserted; $L$ characters are shared.

\item \textbf{False.} Dijkstra can fail with negative edges even without a negative cycle,
because it finalises distances assuming no later (negative) edge can improve them. Use
Bellman--Ford instead.

\item \textbf{(c) $O(mn)$.} $n-1$ passes, each relaxing all $m$ edges.

\item It tells you the graph contains a \textbf{negative cycle} reachable from the source
(distances should have converged after $n-1$ passes if none existed).

\item \textbf{(b) $O(n^3)$.} Three nested loops over $k, i, j$, each $1$ to $n$.

\item \textbf{True.} A topological order lets a single $O(n+m)$ sweep compute shortest paths
(with $\min$) or longest paths (with $\max$); acyclicity makes longest paths well-defined.

\item The cost contributed by split point $k$ is
$M(i,k) + M(k+1,j) + p_{i-1}\,p_k\,p_j$ (cost of the left block, the right block, and the
final multiply of the two resulting matrices).

\end{enumerate}

\end{document}
